%% For tips on how to write a great abstract, have a look at
%%	-	https://www.cdc.gov/stdconference/2018/How-to-Write-an-Abstract_v4.pdf (presentation, start here)
%%	-	https://users.ece.cmu.edu/~koopman/essays/abstract.html
%%	-	https://search.proquest.com/docview/1417403858
%%  - 	https://www.sciencedirect.com/science/article/pii/S037837821830402X

\begin{abstract}
Traditional autonomous navigation systems often rely on expensive sensors and computationally intensive algorithms, making them difficult to deploy on low-cost robotic platforms. This study proposes a lightweight end-to-end navigation framework that combines geometric computer vision and imitation learning for autonomous path following in structured indoor environments. Using a single RGB camera and ArUco fiducial markers, the system computes a homography matrix to transform camera images into a top-down Bird's-Eye-View (BEV) representation. These processed images are then fed into a lightweight Convolutional Neural Network (CNN) trained through behavioral cloning to directly predict steering commands.

The system was evaluated on four track topologies of increasing difficulty: 45-degree Bézier curves, a 90-degree square path, 180-degree U-turns, and a held-out track that was not seen during training. Performance was measured using cross-track error, heading error, and success rate across forty independent randomized-start trials. The controller successfully completed all trials without leaving the lane. On the unseen track, tracking error increased by only about 1.7 times compared to the trained tracks while maintaining a 100\% success rate, suggesting that the network learned transferable path geometry rather than memorizing specific routes. The complete system runs in real time on a Raspberry Pi, demonstrating that reliable autonomous navigation can be achieved using low-cost hardware and a single-camera setup. These results show that combining homography-based scene transformation with deep learning provides a practical alternative to more expensive SLAM-based navigation approaches in structured indoor environments.

Keywords: Autonomous Navigation, Convolutional Neural Networks, Homography, Behavioral Cloning, Mobile Robots, Computer Vision

\end{abstract}